% @file paper63_galois_darstellungen_de.tex
% @brief Galois-Darstellungen und das Langlands-Programm
% @author Michael Fuhrmann
% @date 2026-03-12
% @project Specialist Mathematics Project, Build 143

\documentclass[a4paper,12pt]{amsart}

\usepackage{amsmath, amssymb, amsthm}
\usepackage[utf8]{inputenc}
\usepackage[T1]{fontenc}
\usepackage[ngerman]{babel}
\usepackage{hyperref}
\usepackage{geometry}
\usepackage{enumitem}
\geometry{margin=2.5cm}

% Theoremumgebungen auf Deutsch
\newtheorem{theorem}{Satz}[section]
\newtheorem{lemma}[theorem]{Lemma}
\newtheorem{proposition}[theorem]{Proposition}
\newtheorem{korollar}[theorem]{Korollar}
\newtheorem{conjecture}[theorem]{Vermutung}
\theoremstyle{definition}
\newtheorem{definition}[theorem]{Definition}
\newtheorem{beispiel}[theorem]{Beispiel}
\theoremstyle{remark}
\newtheorem{bemerkung}[theorem]{Bemerkung}

% Mathematische Abkürzungen
\newcommand{\Q}{\mathbb{Q}}
\newcommand{\Z}{\mathbb{Z}}
\newcommand{\Zl}{\mathbb{Z}_\ell}
\newcommand{\Ql}{\mathbb{Q}_\ell}
\newcommand{\Fp}{\mathbb{F}_p}
\newcommand{\Fq}{\mathbb{F}_q}
\newcommand{\Qbar}{\overline{\mathbb{Q}}}
\newcommand{\GQ}{\mathrm{Gal}(\overline{\mathbb{Q}}/\mathbb{Q})}
\newcommand{\GL}{\mathrm{GL}}
\newcommand{\SL}{\mathrm{SL}}
\newcommand{\Frob}{\mathrm{Frob}}
\newcommand{\Spur}{\mathrm{tr}}
\newcommand{\BdR}{B_{\mathrm{dR}}}
\newcommand{\Bcrys}{B_{\mathrm{crys}}}
\newcommand{\Tl}{T_\ell}
\newcommand{\Vl}{V_\ell}

\title{Galois-Darstellungen und das Langlands-Programm}

\author{Michael Fuhrmann}
\address{Specialist Mathematics Project, Build 143}
\email{specialist-maths@project.local}
\date{2026-03-12}

\begin{abstract}
Galois-Darstellungen sind stetige Gruppenhomomorphismen
$\varrho\colon \mathrm{Gal}(\overline{\mathbb{Q}}/\mathbb{Q}) \to \GL_n(R)$,
wobei $R$ typischerweise ein Ring von $\ell$-adischen ganzen Zahlen oder ein endlicher
Körper ist. Sie bilden die zentrale Brücke zwischen algebraischer Zahlentheorie und
analytischen Objekten wie Modulformen und automorphen Darstellungen und sind die
Grundsprache des Langlands-Programms.

Diese Arbeit gibt einen Überblick über die Theorie der Galois-Darstellungen, von der
absoluten Galois-Gruppe über $\ell$-adische und $p$-adische Darstellungen bis hin zu
ihren tiefen Verbindungen mit elliptischen Kurven, Modulformen und der Langlands-Korrespondenz.
Wir behandeln den Modularitätssatz von Wiles--Taylor--Wiles (1995), der bewiesen ist
und Fermats Letzten Satz impliziert; Serres Modularitätsvermutung, bewiesen von
Khare--Wintenberger (2009); Fontaines $p$-adische Hodge-Theorie; sowie Lafforgues
Beweis der Langlands-Korrespondenz über Funktionenkörpern (2002, Fields-Medaille).
Die globale Langlands-Korrespondenz für $\mathbb{Q}$ bleibt eine zentrale offene
Vermutung der Mathematik.
\end{abstract}

\maketitle

\tableofcontents

%=========================================================
\section{Einleitung}
%=========================================================

Die Theorie der Galois-Darstellungen steht am Schnittpunkt zweier großer Strömungen
der modernen Mathematik: der \emph{algebraischen} (Galois-Theorie, algebraische
Zahlkörper, elliptische Kurven, Motive) und der \emph{analytischen} (automorphe
Formen, $L$-Funktionen, harmonische Analysis auf Lie-Gruppen). Eine Galois-Darstellung
codiert die Wirkung der absoluten Galois-Gruppe $\GQ$ auf einem endlichdimensionalen
Vektorraum und übersetzt die von Natur aus unendliche, topologische Gruppe in
konkrete lineare Algebra.

Die leitende Philosophie des \textbf{Langlands-Programms}, das Robert Langlands 1967
und 1970 formulierte, ist das Prinzip der \emph{Funktorialität}: $n$-dimensionale
Galois-Darstellungen sollen bijektiv den automorphen Darstellungen von
$\GL_n(\mathbb{A}_\Q)$, der adelischen allgemeinen linearen Gruppe, entsprechen.
Diese Korrespondenz würde, wenn sie vollständig bewiesen wäre, weite Gebiete der
Mathematik vereinheitlichen und beispiellosen Zugang zu $L$-Funktionen bieten.

\subsection*{Historischer Kontext}

Die Wurzeln der Galois-Darstellungen liegen in klassischen Reziprozitätsgesetzen.
Das quadratische Reziprozitätsgesetz von Gauß (1796) lässt sich als Aussage über das
Verhalten von Primzahlen in quadratischen Körpern reformulieren, kontrolliert durch
eine eindimensionale Galois-Darstellung (ein Dirichlet-Charakter). Die Klassenkörpertheorie,
entwickelt von Hilbert, Takagi und Artin im ersten Drittel des 20.\ Jahrhunderts,
klassifiziert alle abelschen Erweiterungen von Zahlkörpern und begründet das abelsche
Langlands-Programm in Gestalt des Artinschen Reziprozitätsgesetzes.

Das nicht-abelsche Programm begann mit Eichler--Shimura (1950er Jahre), die zeigten,
dass die $L$-Funktion einer Hecke-Eigenform vom Gewicht~2 gleich der $L$-Funktion
einer elliptischen Kurve ist. Serre und Deligne entwickelten in den 1960er und 1970er
Jahren die Theorie der $\ell$-adischen Darstellungen systematisch. Der spektakuläre
Höhepunkt war Wiles' Beweis des Letzten Fermatschen Satzes (1995), der auf dem
Modularitätssatz für semistabile elliptische Kurven beruht.

\subsection*{Aufbau dieser Arbeit}

Abschnitte~\ref{sec:galois-gruppe} und~\ref{sec:ell-adisch} legen die absolute
Galois-Gruppe und $\ell$-adische Darstellungen fest. Abschnitt~\ref{sec:elliptische-kurven}
konstruiert Darstellungen aus elliptischen Kurven über den Tate-Modul.
Abschnitte~\ref{sec:modularitaet} und~\ref{sec:modulformen} behandeln den Modularitätssatz
und modulare Galois-Darstellungen. Abschnitte~\ref{sec:serre} und~\ref{sec:fontaine}
befassen sich mit Serres Modularitätssatz und Fontaines $p$-adischer Hodge-Theorie.
Abschnitte~\ref{sec:langlands} und~\ref{sec:neuere-resultate} geben einen Überblick
über die Langlands-Korrespondenz und neuere Ergebnisse.

%=========================================================
\section{Die absolute Galois-Gruppe}
\label{sec:galois-gruppe}
%=========================================================

\begin{definition}
Die \textbf{absolute Galois-Gruppe} von $\Q$ ist
\[
  G_\Q := \GQ = \mathrm{Aut}(\Qbar/\Q),
\]
die Gruppe aller Körperautomorphismen eines festen algebraischen Abschlusses $\Qbar$,
die $\Q$ punktweise fixieren.
\end{definition}

\subsection{Profinite Topologie}

Obwohl $G_\Q$ riesig (überabzählbar) ist, trägt sie eine natürliche Topologie, die
sie zu einer profiniten Gruppe macht:
\[
  G_\Q \;\cong\; \varprojlim_{[K:\Q] < \infty,\; K/\Q \text{ galoissch}} \mathrm{Gal}(K/\Q).
\]
Hierbei ist der inverse Limes über das gerichtete System endlicher Galois-Erweiterungen
$K/\Q$, geordnet durch Inklusion. Die Topologie ist die \emph{Krull-Topologie}: eine
Basis offener Untergruppen besteht aus den Normaluntergruppen $\mathrm{Gal}(\Qbar/K)$
für endliche galoissche $K/\Q$.

\begin{proposition}
$G_\Q$ ist kompakt, total unzusammenhängend und hausdorffsch (d.\,h.\ profinit).
\end{proposition}

\begin{proof}
Nach dem Satz von Tychonoff ist das Produkt $\prod_K \mathrm{Gal}(K/\Q)$ über alle
endlichen Galois-Erweiterungen kompakt. Der inverse Limes ist eine abgeschlossene
Untergruppe, also ebenfalls kompakt. Totale Unzusammenhängtheit folgt daraus, dass
jede offene Untergruppe endlichen Index hat und die offenen Untergruppen eine
Umgebungsbasis bilden.
\end{proof}

\subsection{Frobenius-Elemente}

Für eine bei $p$ unverzweigte Galois-Erweiterung $K/\Q$ gibt es eine wohldefinierte
Konjugiertenklasse $\Frob_p \in \mathrm{Gal}(K/\Q)$, charakterisiert durch
\[
  \Frob_p(\alpha) \equiv \alpha^p \pmod{\mathfrak{P}}
\]
für alle algebraischen ganzen Zahlen $\alpha \in \mathcal{O}_K$ und jedes Primideal
$\mathfrak{P}$ von $K$ über $p$.

\begin{definition}
Für eine Primzahl $\ell$ ist das \textbf{geometrische Frobenius-Element} an $p \neq \ell$
das Element $\Frob_p \in G_\Q$ (bzw.\ die von ihm definierte Konjugiertenklasse),
das auf den $\ell$-Potenz-Einheitswurzeln durch $\zeta_{\ell^n} \mapsto \zeta_{\ell^n}^p$
wirkt.
\end{definition}

Nach dem Dichtigkeitssatz von Chebotarev sind die Frobenius-Elemente
$\{\Frob_p : p \text{ unverzweigt}\}$ in $G_\Q$ \emph{dicht}. Folglich ist eine
stetige Darstellung vollständig durch die Werte an Frobenius-Elementen bestimmt.

\subsection{Verzweigung}

Für jede Primzahl $p$ fixiere eine Einbettung $\Qbar \hookrightarrow \overline{\Q}_p$.
Dies ergibt eine Zerlegungsgruppe $D_p \subset G_\Q$ und eine Trägheitsgruppe
$I_p \trianglelefteq D_p$, eingebettet in die exakte Folge
\[
  1 \to I_p \to D_p \to G_{\Fp} \cong \hat{\Z} \to 1.
\]
Das arithmetische Frobenius-Erzeugende von $G_{\Fp}$ wird zu $\Frob_p \in D_p/I_p$ angehoben.

%=========================================================
\section{\texorpdfstring{$\ell$}{l}-adische Darstellungen}
\label{sec:ell-adisch}
%=========================================================

Sei $\ell$ eine Primzahl. Der Ring $\Zl$ der $\ell$-adischen ganzen Zahlen ist die
Vervollständigung von $\Z$ bei $\ell$, und $\Ql = \Zl[\ell^{-1}]$ ist der Körper der
$\ell$-adischen Zahlen.

\begin{definition}
Eine \textbf{$\ell$-adische Darstellung} von $G_\Q$ der Dimension $n$ ist ein stetiger
Homomorphismus
\[
  \varrho\colon G_\Q \to \GL_n(\Zl),
\]
wobei $\GL_n(\Zl)$ die $\ell$-adische Topologie trägt. Oft betrachtet man auch die
zugehörige rationale Darstellung $\varrho\colon G_\Q \to \GL_n(\Ql)$.
\end{definition}

\subsection{Stetigkeit und Endlichkeit}

Da $G_\Q$ profinit und $\GL_n(\Zl)$ eine $\ell$-adische Lie-Gruppe ist, bedeutet
Stetigkeit, dass der Kern eine offene Untergruppe enthält, d.\,h.\ dass $\varrho$
durch eine offene Untergruppe endlichen Indexes von $G_\Q$ faktorisiert.

\subsection{Der zyklotomische Charakter}

Die grundlegende eindimensionale $\ell$-adische Darstellung ist der
\textbf{zyklotomische Charakter}
\[
  \chi_\ell\colon G_\Q \to \Zl^\times,
\]
definiert durch $\sigma(\zeta) = \zeta^{\chi_\ell(\sigma)}$ für alle $\ell^n$-ten
Einheitswurzeln $\zeta$. Konkret gilt via den kanonischen Isomorphismus
$\hat{\Z} \cong \varprojlim \Z/\ell^n\Z$:
\[
  \chi_\ell(\Frob_p) = p \quad \text{für alle Primzahlen } p \neq \ell.
\]

\begin{beispiel}[Dirichlet-Charaktere]
Ein Dirichlet-Charakter $\chi\colon (\Z/N\Z)^\times \to \mathbb{C}^\times$ des
Führers $N$ liefert über die Klassenkörpertheorie eine eindimensionale Galois-Darstellung:
\[
  \varrho_\chi\colon G_\Q \twoheadrightarrow \mathrm{Gal}(\Q(\zeta_N)/\Q)
  \cong (\Z/N\Z)^\times \xrightarrow{\chi} \mathbb{C}^\times.
\]
Der Frobenius-Wert ist $\varrho_\chi(\Frob_p) = \chi(p)$ für $p \nmid N$.
\end{beispiel}

\subsection{Tate-Drehungen}

Für $n \in \Z$ ist die $n$-fache \textbf{Tate-Drehung} einer Darstellung $V$ durch
$V(n) := V \otimes \chi_\ell^{\otimes n}$ definiert. Die Kategorie der $\ell$-adischen
Darstellungen ist abgeschlossen unter direkten Summen, Tensorprodukten, Dualen und
Tate-Drehungen.

%=========================================================
\section{Darstellungen aus elliptischen Kurven}
\label{sec:elliptische-kurven}
%=========================================================

Elliptische Kurven liefern die wichtigsten zweidimensionalen Galois-Darstellungen.

\begin{definition}
Sei $E/\Q$ eine elliptische Kurve und $\ell$ eine Primzahl. Der
\textbf{$\ell$-adische Tate-Modul} ist
\[
  \Tl(E) := \varprojlim_n E[\ell^n],
\]
wobei $E[\ell^n] \cong (\Z/\ell^n\Z)^2$ die $\ell^n$-Torsionspunkte (über $\Qbar$)
sind und die Übergangsabbildungen Multiplikation mit $\ell$. Als abelsche Gruppe gilt
$\Tl(E) \cong \Zl^2$.
\end{definition}

Da $G_\Q$ die Torsionspunkte permutiert, wirkt sie auf $\Tl(E)$ und liefert:

\begin{definition}
Die \textbf{$\ell$-adische Darstellung von $E$} ist
\[
  \varrho_{E,\ell}\colon G_\Q \to \GL_2(\Zl).
\]
Ihre rationale Version $V_\ell(E) := \Tl(E) \otimes_{\Zl} \Ql$ liefert
$\varrho_{E,\ell}\colon G_\Q \to \GL_2(\Ql)$.
\end{definition}

\subsection{Eigenschaften von $\varrho_{E,\ell}$}

\begin{theorem}[Weil 1940]
\label{satz:weil-hasse}
Sei $E/\Q$ eine elliptische Kurve mit guter Reduktion bei $p \neq \ell$.
Dann ist $\varrho_{E,\ell}$ bei $p$ unverzweigt, und
\[
  \det(1 - \varrho_{E,\ell}(\Frob_p)\,X) = 1 - a_p X + p X^2,
\]
wobei $a_p = p + 1 - \#E(\Fp)$. Insbesondere gilt
$|\alpha_p|, |\beta_p| = \sqrt{p}$ (Hassescher Satz: $|a_p| \leq 2\sqrt{p}$).
\end{theorem}

\begin{proof}[Beweisskizze]
Das charakteristische Polynom von $\Frob_p$ auf $V_\ell(E)$ berechnet die
Zeta-Funktion von $E$ über $\Fp$ via der Lefschetz-Spurformel am Tate-Modul.
Hasses Ungleichung $|a_p| \leq 2\sqrt{p}$ folgt aus der Positivität der
Rosati-Involution.
\end{proof}

\subsection{Verzweigung an Primzahlen schlechter Reduktion}

An einer Primzahl $p$ additiver Reduktion ist $\varrho_{E,\ell}|_{I_p}$ nicht-trivial
und erfasst die lokale Néron-Modell-Struktur. Bei semistabiler Reduktion
(multiplikative Reduktion) wirkt die Trägheit unipotent (Tate-Kurven-Theorie).

\subsection{Das Weil-Paarung}

Die natürliche alternierende Paarung $e_\ell\colon \Tl(E) \times \Tl(E) \to \Tl(\mu)$
liefert
\[
  \det(\varrho_{E,\ell}) = \chi_\ell,
\]
den zyklotomischen Charakter. Dies ist die Schlüsselbeziehung, die die
zweidimensionale Darstellung an den eindimensionalen zyklotomischen Charakter koppelt.

%=========================================================
\section{Der Modularitätssatz}
\label{sec:modularitaet}
%=========================================================

\begin{theorem}[Modularitätssatz; Wiles 1995, Taylor--Wiles 1995,
  Breuil--Conrad--Diamond--Taylor 2001]
\label{satz:modularitaet}
Jede elliptische Kurve $E$ über $\Q$ ist modular: Es existiert eine Neuform
$f = \sum_{n=1}^\infty a_n(f) q^n$ vom Gewicht~2 und einem gewissen Führer $N$, sodass
\[
  a_p(E) = a_p(f) \quad \text{für alle Primzahlen } p \nmid N.
\]
Äquivalent: $L(E, s) = L(f, s)$.
\end{theorem}

\begin{bemerkung}
Wiles bewies den semistabilen Fall (1995), der für den Beweis von Fermats Letztem
Satz ausreichte. Der vollständige Satz für alle elliptischen Kurven über $\Q$ wurde
von Breuil, Conrad, Diamond und Taylor (2001) bewiesen.
\end{bemerkung}

\subsection{Beweisidee}

Der Beweis verläuft in folgenden Schritten:

\begin{enumerate}[label=(\roman*)]
\item \textbf{Deformationstheorie}: Für eine Primzahl $p \mid N$ (Wiles wählt $p = 3$
oder $p = 5$) betrachte die mod-$p$-Darstellung $\bar{\varrho}_{E,p}$.

\item \textbf{Modularitätsliftung}: Zeige, dass $\varrho_{E,p}$ modular ist, durch
Liften der residuellen mod-$p$-Darstellung $\bar{\varrho}_{E,p}$ mit dem
\emph{numerischen Kriterium} (Taylor--Wiles-Methode).

\item \textbf{Taylor--Wiles-Methode}: Konstruiere ein Hilfssystem von
Galois-Darstellungen (die \emph{Hecke-Algebren}) und beweise einen $R = T$-Satz
(Deformationsring $R$ gleich Hecke-Algebra $T$) durch einen Dimensionsvergleich.

\item \textbf{Basisfall}: Nutze die bekannte Modularität von $E$ modulo kleiner
Primzahlen über Langlands--Tunnell (für $p = 3$).
\end{enumerate}

\subsection{Fermats Letzter Satz als Korollar}

\begin{korollar}[Fermats Letzter Satz; Wiles 1995]
Für $n \geq 3$ hat die Gleichung $x^n + y^n = z^n$ keine Lösungen in ganzen
Zahlen mit $xyz \neq 0$.
\end{korollar}

\begin{proof}[Beweis (unter Verwendung von Satz~\ref{satz:modularitaet})]
Angenommen, $a^p + b^p = c^p$ mit $abc \neq 0$ und Primzahl $p \geq 5$ (kleine Fälle
sind klassisch). Frey (1986) konstruierte die elliptische Kurve
$E_{a,b,c}\colon y^2 = x(x-a^p)(x+b^p)$. Ribet (1990) bewies (unter Verwendung des
Mazur-Prinzips), dass wenn $E_{a,b,c}$ modular wäre, die zugehörige Neuform Führer
$N = 2$ hätte --- aber es gibt keine Neuformen vom Gewicht~2 und Führer~2 (der
betreffende Raum ist 0-dimensional). Dies widerspricht der Modularität, also existiert
keine Frey-Kurve, also keine Fermat-Lösung.
\end{proof}

%=========================================================
\section{Modulare Galois-Darstellungen}
\label{sec:modulformen}
%=========================================================

Sei $f = \sum_{n \geq 1} a_n q^n$ eine normalisierte Hecke-Eigenform vom Gewicht~$k$,
Führer~$N$ und Nebentypus~$\chi$. Sei $K_f = \Q(\{a_n\})$ ihr Koeffizientenkörper.

\begin{theorem}[Deligne 1971]
\label{satz:deligne}
Für jede Primzahl $\ell$ und Einbettung $\iota\colon K_f \hookrightarrow \overline{\Q}_\ell$
existiert eine eindeutige (bis auf Isomorphie) irreduzible, stetige Darstellung
\[
  \varrho_{f,\ell}\colon G_\Q \to \GL_2(\mathcal{O}_{K_f \otimes \Ql}),
\]
unverzweigt außerhalb $N\ell$, so dass für alle Primzahlen $p \nmid N\ell$:
\[
  \Spur(\varrho_{f,\ell}(\Frob_p)) = a_p(f), \qquad
  \det(\varrho_{f,\ell}(\Frob_p)) = \chi(p)\, p^{k-1}.
\]
\end{theorem}

\subsection{Die Ramanujan--Deligne-Schranke}

\begin{theorem}[Deligne 1974; Weil-Vermutungen für Varietäten]
\label{satz:ramanujan-deligne}
Für eine Neuform $f$ vom Gewicht~$k$ wie oben und $p \nmid N$ erfüllen die Eigenwerte
$\alpha_p, \beta_p$ von $\varrho_{f,\ell}(\Frob_p)$:
\[
  |\iota(\alpha_p)| = |\iota(\beta_p)| = p^{(k-1)/2}
\]
für jede Einbettung $\iota\colon K_f \hookrightarrow \mathbb{C}$.
Für $k = 2$ ergibt sich $|a_p| \leq 2\sqrt{p}$ (Ramanujans Vermutung für $\Delta$
und ihre Verallgemeinerung).
\end{theorem}

\begin{proof}[Beweisidee]
Deligne realisiert $\varrho_{f,\ell}$ in der étalen Kohomologie $H^{k-1}_{\text{ét}}$
einer Kuga--Sato-Varietät und wendet seinen Beweis der Weil-Vermutungen (Riemannsche
Hypothese für Varietäten über endlichen Körpern) an, der die Schranke $p^{(k-1)/2}$
an die Frobenius-Eigenwerte liefert.
\end{proof}

%=========================================================
\section{Serres Modularitätssatz}
\label{sec:serre}
%=========================================================

\begin{definition}
Eine \textbf{mod-$\ell$-Galois-Darstellung} ist ein stetiger Homomorphismus
$\bar{\varrho}\colon G_\Q \to \GL_2(\overline{\Fp})$.
Sie heißt \textbf{ungerade}, wenn $\det(\bar{\varrho}(\text{komplexe Konjugation})) = -1$,
und \textbf{irreduzibel}, wenn sie keine $G_\Q$-stabile Gerade besitzt.
\end{definition}

\begin{conjecture}[Serre 1987]
Jede ungerade, irreduzible stetige Darstellung
$\bar{\varrho}\colon G_\Q \to \GL_2(\overline{\mathbb{F}}_\ell)$
stammt von einer Neuform vom Gewicht $k(\bar{\varrho})$ und Führer $N(\bar{\varrho})$,
wie durch Serres Rezept vorgeschrieben.
\end{conjecture}

\begin{theorem}[Khare--Wintenberger 2009; Serres Modularitätssatz]
\label{satz:serre}
Serres Vermutung ist wahr: Jede ungerade irreduzible
$\bar{\varrho}\colon G_\Q \to \GL_2(\overline{\mathbb{F}}_\ell)$ ist modular.
\end{theorem}

\begin{proof}[Beweisidee]
Khare und Wintenberger beweisen dies durch eine zweistu\-fige Induktion:
\begin{enumerate}[label=(\alph*)]
\item \textbf{Niveauabsenkung}: Verwendung von Ribets Methode und Modularitätsliftungs\-sätzen
  (im Geiste von Wiles), um auf den Fall kleinen Niveaus und Gewichts zu reduzieren.
\item \textbf{Gewichtsabsenkung}: Reduktion auf charakteristik-$p$-Darstellungen
  kleinen Führers mittels $p$-adischer Hodge-Theorie und Fontaine--Mazur-artiger Argumente.
\item \textbf{Basisfall}: Überprüfung der Vermutung für kleines Gewicht und Niveau
  durch explizite endliche Berechnungen.
\end{enumerate}
Der Beweis verwendet die volle Stärke der Taylor--Wiles--Kisin-Modularitätsliftungs\-maschinerie.
\end{proof}

\begin{korollar}
Serres Modularitätssatz impliziert den Modularitätssatz (Satz~\ref{satz:modularitaet})
und vertieft das Verständnis, wie residuelle Darstellungen zu Darstellungen der
Charakteristik Null angehoben werden.
\end{korollar}

%=========================================================
\section{\texorpdfstring{$p$}{p}-adische Darstellungen und Fontaines Theorie}
\label{sec:fontaine}
%=========================================================

Die lokale Theorie der $p$-adischen Darstellungen (d.\,h.\ stetige
$G_{\Q_p}$-Darstellungen auf $\Q_p$-Vektorräumen) wird durch Fontaines
\emph{$p$-adische Hodge-Theorie} kontrolliert.

\subsection{Periodenringe}

Fontaine konstruierte mehrere Periodenringe zur Klassifikation $p$-adischer Darstellungen:

\begin{definition}
\begin{itemize}
  \item $\BdR$: der de-Rham-Periodenring; ein vollständiger diskreter Bewertungskörper
    mit Restklassenkörper $\mathbb{C}_p$ und einer Filtration.
  \item $\Bcrys \subset \BdR$: der kristalline Periodenring mit Frobenius $\phi$
    und einer Filtration.
  \item $B_{\mathrm{st}} \supset \Bcrys$: der semistabile Periodenring mit
    zusätzlichem Monodromie-Operator $N$.
\end{itemize}
\end{definition}

\begin{definition}
Eine $p$-adische Darstellung $V$ von $G_{\Q_p}$ heißt:
\begin{itemize}
  \item \textbf{Hodge--Tate}, wenn $V \otimes_{\Q_p} \mathbb{C}_p \cong \bigoplus_i \mathbb{C}_p(i)^{h_i}$,
  \item \textbf{de Rham}, wenn $D_{\mathrm{dR}}(V) := (V \otimes_{\Q_p} \BdR)^{G_{\Q_p}}$ die Dimension $n = \dim V$ hat,
  \item \textbf{kristallin}, wenn $D_{\mathrm{crys}}(V) := (V \otimes_{\Q_p} \Bcrys)^{G_{\Q_p}}$ die Dimension $n$ hat,
  \item \textbf{semistabil}, wenn $D_{\mathrm{st}}(V) := (V \otimes_{\Q_p} B_{\mathrm{st}})^{G_{\Q_p}}$ die Dimension $n$ hat.
\end{itemize}
\end{definition}

\begin{theorem}[Fontaine--Faltings--Tsuji]
\label{satz:fontaine-vergleich}
Die $\ell$-adische Darstellung $V_\ell(E)$ einer elliptischen Kurve $E/\Q_p$ ist:
\begin{itemize}
  \item Kristallin, wenn $E$ gute Reduktion bei $p$ hat,
  \item Semistabil, wenn $E$ semistabile (multiplikative) Reduktion bei $p$ hat.
\end{itemize}
\end{theorem}

\subsection{Wach-Moduln und ganzzahlige Strukturen}

Für kristalline Darstellungen liefern Wach-Moduln ganzzahlige ($\Zl$-Gitter-)
Versionen der Fontaine-Klassifikation, die das Studium von Darstellungen über $\Zl$
statt $\Ql$ ermöglichen.

\subsection{Die Fontaine--Mazur-Vermutung}

\begin{conjecture}[Fontaine--Mazur 1993]
Eine irreduzible $p$-adische Darstellung $\varrho\colon G_\Q \to \GL_n(\Ql)$
stammt aus der Geometrie (d.\,h.\ aus der étalen Kohomologie einer Varietät über~$\Q$)
genau dann, wenn sie bei $p$ \emph{de Rham} ist (lokal bei $p$) und bei fast allen
Primzahlen unverzweigt ist.
\end{conjecture}

Diese Vermutung ist in wichtigen Spezialfällen bewiesen (z.\,B.\ $n = 2$, in
Verbindung mit Modularitätsliftung).

%=========================================================
\section{Die Langlands-Korrespondenz}
\label{sec:langlands}
%=========================================================

\subsection{Die globale Langlands-Korrespondenz für $\GL_n$}

\begin{conjecture}[Globale Langlands-Korrespondenz für $\Q$]
\label{verm:langlands}
Es gibt eine Bijektion zwischen:
\begin{enumerate}[label=(\roman*)]
\item $n$-dimensionalen irreduziblen stetigen Darstellungen
  $\varrho\colon G_\Q \to \GL_n(\overline{\Q}_\ell)$ (geometrischer Herkunft),
\item kuspidalen automorphen Darstellungen $\pi$ von $\GL_n(\mathbb{A}_\Q)$,
\end{enumerate}
dergestalt dass $L(s, \varrho) = L(s, \pi)$ für die zugehörigen $L$-Funktionen.
\end{conjecture}

Diese Vermutung bleibt für $\Q$ im Allgemeinen \textbf{offen}. Jedoch sind
Spezialfälle etabliert:
\begin{itemize}
\item $n = 1$: Klassenkörpertheorie / Artinsches Reziprozitätsgesetz (bewiesen).
\item $n = 2$: Modulformen vom Gewicht~$k$ (für elliptische Kurven durch Wiles et al.
  bewiesen, allgemeiner durch Deligne für modulare Galois-Darstellungen).
\item Funktionenkörper: Bewiesen durch Lafforgue (2002), s.\ Abschnitt~\ref{sec:neuere-resultate}.
\end{itemize}

\subsection{Lokale Langlands-Korrespondenz}

\begin{theorem}[Lokale Langlands-Korrespondenz; Langlands 1970, Henniart 2000, Harris--Taylor 2001]
\label{satz:lokale-langlands}
Für jede Primzahl $p$ gibt es eine natürliche Bijektion
\[
  \left\{
    \begin{array}{c}
      n\text{-dim.\ Frobenius-semi-}\\ \text{einfache } W_{D}(\Q_p)\text{-Darst.}
    \end{array}
  \right\}
  \xrightarrow{\;\sim\;}
  \left\{
    \begin{array}{c}
      \text{irred.\ glatte Darst.}\\ \text{von } \GL_n(\Q_p)
    \end{array}
  \right\},
\]
wobei $W_D(\Q_p)$ die Weil--Deligne-Gruppe bezeichnet.
\end{theorem}

\subsection{Funktorialität}

Langlands' Funktorialitätsprinzip besagt, dass ein Homomorphismus von $L$-Gruppen
${}^L G \to {}^L H$ einem Transfer automorpher Darstellungen von $G(\mathbb{A})$
nach $H(\mathbb{A})$ entsprechen soll. In vielen Fällen bewiesen (Basiswechsel,
symmetrische Potenzhebungen für $\GL_2$), im Allgemeinen aber noch offen.

%=========================================================
\section{Neuere Resultate}
\label{sec:neuere-resultate}
%=========================================================

\subsection{Lafforgues Satz über Funktionenkörpern}

\begin{theorem}[Lafforgue 2002; Fields-Medaille]
\label{satz:lafforgue}
Die globale Langlands-Korrespondenz gilt für $\GL_n$ über einem Funktionenkörper
$\mathbb{F}_q(C)$ (dem Funktionenkörper einer Kurve $C$ über einem endlichen Körper
$\mathbb{F}_q$): Es gibt eine Bijektion zwischen $n$-dimensionalen $\ell$-adischen
Darstellungen der absoluten Galois-Gruppe von $\mathbb{F}_q(C)$ und kuspidalen
automorphen Darstellungen von $\GL_n(\mathbb{A}_{\mathbb{F}_q(C)})$, die
$L$-Funktionen in Übereinstimmung bringt.
\end{theorem}

Dies war der erste vollständige Beweis der globalen Langlands-Korrespondenz für
$\GL_n$ mit $n > 2$ und beruht auf Drinfelds Arbeit für $\GL_2$ (1974, 1980).

\subsection{Potenzielle Automorphie}

\begin{theorem}[Taylor 2004; Harris--Shepherd-Barron--Taylor 2010]
\label{satz:potenzielle-automorphie}
Sei $\varrho\colon G_\Q \to \GL_n(\overline{\Q}_\ell)$ eine „reguläre, ungerade"
$\ell$-adische Darstellung. Dann ist $\varrho$ \emph{potenziell automorph}: Es
existiert ein total reeller Körper $F$, sodass $\varrho|_{G_F}$ automorph ist.
\end{theorem}

Potenzielle Automorphie genügt, um zu beweisen, dass $L(\varrho, s)$ analytische
Fortsetzung besitzt und eine Funktionalgleichung erfüllt, auch wenn volle Automorphie
nicht etabliert ist.

\subsection{Die Taylor--Wiles--Kisin-Methode}

Diamond (1997), Kisin (2009) und viele andere erweiterten die Taylor--Wiles-Methode
auf alle zweidimensionalen de-Rham-Darstellungen von $G_{\Q_p}$. Kisins Beweis
verwendet \emph{Breuil--Kisin-Moduln} (Analoga der Wach-Moduln) und löst die meisten
Fälle der Fontaine--Mazur-Vermutung für $n = 2$.

\subsection{$p$-adische Langlands-Korrespondenz}

Für $\GL_2(\Q_p)$ etablierte Colmez (2010) eine $p$-adische lokale Langlands-Korrespondenz:
eine Bijektion zwischen zweidimensionalen triangulinären $p$-adischen Darstellungen
von $G_{\Q_p}$ und lokal-analytischen Darstellungen von $\GL_2(\Q_p)$, verträglich
mit der klassischen lokalen Langlands-Korrespondenz.

\subsection{Höherdimensionale und geometrische Fortschritte}

Neuere Arbeiten von Scholze (Verwendung von perfektoiden Räumen und Diamanten),
V.~Lafforgue (2018, für reduktive Gruppen über Funktionenkörpern) und
Fargues--Scholze (2021, Geometrisierung des lokalen Langlands-Programms) haben das
strukturelle Verständnis drastisch erweitert. Fargues und Scholze haben das lokale
Langlands-Programm geometrisiert und eine garbentheoretische Version über der
Fargues--Fontaine-Kurve konstruiert.

%=========================================================
\section{Literatur}
%=========================================================

\begin{thebibliography}{99}

\bibitem{Wiles1995}
A.~Wiles,
\textit{Modular elliptic curves and Fermat's last theorem},
Annals of Mathematics \textbf{141} (1995), 443--551.

\bibitem{TaylorWiles1995}
R.~Taylor und A.~Wiles,
\textit{Ring-theoretic properties of certain Hecke algebras},
Annals of Mathematics \textbf{141} (1995), 553--572.

\bibitem{BCDT2001}
C.~Breuil, B.~Conrad, F.~Diamond und R.~Taylor,
\textit{On the modularity of elliptic curves over $\mathbb{Q}$: Wild 3-adic exercises},
Journal of the American Mathematical Society \textbf{14} (2001), 843--939.

\bibitem{KhareWintenberger2009}
C.~Khare und J.-P.~Wintenberger,
\textit{Serre's modularity conjecture (I)},
Inventiones Mathematicae \textbf{178} (2009), 485--504.

\bibitem{Serre1987}
J.-P.~Serre,
\textit{Sur les représentations modulaires de degré 2 de $\mathrm{Gal}(\overline{\mathbb{Q}}/\mathbb{Q})$},
Duke Mathematical Journal \textbf{54} (1987), 179--230.

\bibitem{Fontaine1994}
J.-M.~Fontaine,
\textit{Représentations $p$-adiques semi-stables},
Astérisque \textbf{223} (1994), 113--184.

\bibitem{Lafforgue2002}
L.~Lafforgue,
\textit{Chtoucas de Drinfeld et correspondance de Langlands},
Inventiones Mathematicae \textbf{147} (2002), 1--241.

\bibitem{Deligne1974}
P.~Deligne,
\textit{La conjecture de Weil. I},
Publications Mathématiques de l'IHÉS \textbf{43} (1974), 273--307.

\bibitem{Langlands1970}
R.~P.~Langlands,
\textit{Problems in the theory of automorphic forms},
Lectures in Modern Analysis and Applications III,
Lecture Notes in Mathematics \textbf{170}, Springer, 1970, 18--61.

\bibitem{HarrisTaylor2001}
M.~Harris und R.~Taylor,
\textit{The geometry and cohomology of some simple Shimura varieties},
Annals of Mathematics Studies \textbf{151}, Princeton University Press, 2001.

\end{thebibliography}

\end{document}
